% =====================================================================
%  Figura 4.2 — Biomarcadores para Longevidade: Faixas Normais vs. Ótimas
%  Estrutura igual ao original colorido: 4 seções separadas, cada uma com
%  cabeçalho próprio + headers de coluna repetidos. Coluna "Faixa ótima"
%  com background cinza claro (substitui o gold/cream do colorido).
% =====================================================================
\begingroup
\sffamily

% Título principal
\begin{center}
{\sffamily\bfseries\color{petrol}\normalsize%
Biomarcadores para longevidade: faixas normais vs.\ ótimas\par}
{\sffamily\scriptsize\color{ink}%
16 biomarcadores principais + 2 complementares, organizados em 4 grupos temáticos.\par}
\end{center}

\vspace{0.2em}

% Definições de coluna comuns — compactadas para caber em 1 página
\newcolumntype{B}{>{\raggedright\arraybackslash}p{0.18\textwidth}}
\newcolumntype{N}{>{\raggedright\arraybackslash}p{0.20\textwidth}}
\newcolumntype{O}{>{\columncolor{gray!12}\raggedright\arraybackslash}p{0.22\textwidth}}
\newcolumntype{R}{>{\raggedright\arraybackslash}p{0.32\textwidth}}

\renewcommand{\arraystretch}{1.02}
\setlength{\tabcolsep}{3pt}

% Macro para o cabeçalho de uma seção (4 column headers — repetidos a cada bloco)
\newcommand{\biotabhead}{%
\rowcolor{gray!22}%
{\tiny\textbf{BIOMARCADOR}} &
{\tiny\textbf{FAIXA "NORMAL" (LAB)}} &
{\tiny\textbf{FAIXA ÓTIMA (LONGEVIDADE)}} &
{\tiny\textbf{O QUE REVELA}} \\
\midrule
}

% Macro para o título da seção
\newcommand{\biotabsection}[2]{%
\vspace{0.25em}\par%
{\sffamily\bfseries\color{petrol}\small #1.\ #2}\par%
\vspace{0.1em}%
}

% ===================== Seção 1: Metabólico =====================
\biotabsection{1}{METABÓLICO}
{\tiny\noindent
\begin{tabular}{@{}B N O R@{}}
\toprule
\biotabhead
Insulina de jejum & < 25 µUI/mL & \textbf{< 8 µUI/mL} & Resistência insulínica precoce — visível 5--10 anos antes da glicose \\
HbA1c & < 6{,}5\% & \textbf{< 5{,}4\% (ideal 4{,}8--5{,}2\%)} & Velocidade do envelhecimento glicêmico — média de 2--3 meses \\
Relação TG/HDL & sem ref.\ clínica & \textbf{< 2{,}0 (> 3{,}5 = risco)} & Proxy acessível para resistência insulínica \\
\bottomrule
\end{tabular}}

% ===================== Seção 2: Inflamação e estresse =====================
\biotabsection{2}{INFLAMAÇÃO E ESTRESSE}
{\tiny\noindent
\begin{tabular}{@{}B N O R@{}}
\toprule
\biotabhead
hs-CRP & < 3{,}0 mg/L & \textbf{< 1{,}0 mg/L} & Inflamação subclínica — base da maioria das doenças (\textit{inflammaging}) \\
Ácido úrico & 3{,}5--7{,}2 mg/dL (M) / 2{,}6--6{,}0 (F) & \textbf{< 5{,}5 mg/dL} & Inflamação metabólica — risco CV independente \\
GGT & < 45 U/L (M) / < 40 (F) & \textbf{< 25--30 U/L} & Estresse oxidativo, esteatose hepática precoce \\
\bottomrule
\end{tabular}}

% ===================== Seção 3: Lipídico e cardiovascular =====================
\biotabsection{3}{LIPÍDICO E CARDIOVASCULAR}
{\tiny\noindent
\begin{tabular}{@{}B N O R@{}}
\toprule
\biotabhead
ApoB & < 130 mg/dL & \textbf{< 90 mg/dL (ideal < 80)} & Marcador causal real — número de partículas aterogênicas \\
Relação ApoB/ApoA1 & sem ref.\ clínica & \textbf{< 0{,}80} & Balanço entre risco e defesa arterial \\
Lp(a) & < 125 nmol/L & \textbf{< 75 nmol/L} & Risco genético cardiovascular — medido uma única vez na vida \\
Colesterol não-HDL & < 130 mg/dL & \textbf{< 100 mg/dL} & Carga lipídica aterogênica em conjunto \\
Troponina I ultrassensível & < 16 ng/L (M) / < 34 (F) & \textbf{indetectável} & Lesão miocárdica subclínica \\
NT-proBNP & < 125 pg/mL & \textbf{< 50 pg/mL (< 50 anos)} & Estresse de parede cardíaca precoce \\
\bottomrule
\end{tabular}}

% ===================== Seção 4: Regulação e longevidade =====================
\biotabsection{4}{REGULAÇÃO E LONGEVIDADE}
{\tiny\noindent
\begin{tabular}{@{}B N O R@{}}
\toprule
\biotabhead
Vitamina D (25-OH) & > 20 ng/mL & \textbf{40--60 ng/mL} & Função imune, óssea e metabólica \\
Homocisteína & < 15 µmol/L & \textbf{< 10 µmol/L} & Risco vascular e cognitivo — reversível com B\textsubscript{12}, B\textsubscript{6}, folato \\
Ferritina & 12--300 ng/mL & \textbf{40--150 ng/mL} & Reserva de ferro \emph{e} proteína inflamatória \\
TFGe (cistatina C) & > 60 mL/min & \textbf{> 90 mL/min} & Função renal precoce \\
Microalbuminúria & < 30 mg/g & \textbf{< 10 mg/g} & Lesão glomerular precoce \\
Albumina sérica & 3{,}5--5{,}0 g/dL & \textbf{$\geq$ 4{,}5 g/dL} & Síntese hepática e estado nutricional \\
\bottomrule
\end{tabular}}

\vspace{0.25em}
{\scriptsize\itshape\color{ink}\noindent%
\textbf{Figura 4.2.}\ Biomarcadores para longevidade: faixas normais vs.\ ótimas.
M = homens; F = mulheres. Faixas de referência podem variar conforme método/laboratório.
Interpretar com painel clínico completo (ESC/EAS, ACC/AHA, ABRAN).\par}
\endgroup

\medskip
